% \omiss produces '[...]'
\newcommand{\omissis}{[\dots\negthinspace]}

% Itemize symbols
% see: https://tex.stackexchange.com/a/62497
% \renewcommand{\labelitemi}{$\bullet$}
% \renewcommand{\labelitemii}{$\cdot$}
% \renewcommand{\labelitemiii}{$\diamond$}
% \renewcommand{\labelitemiv}{$\ast$}

\pagestyle{fancy} 
\setlength{\headheight}{14.5pt}
\fancyhf{} % Svuota le vecchie impostazioni

% Configurazione delle intestazioni: testo non tutto in maiuscolo (nouppercase)
\fancyhead[L]{\fontsize{12}{14.5}\selectfont\nouppercase{\leftmark}}
\fancyhead[R]{\fontsize{12}{14.5}\selectfont\nouppercase{\rightmark}}
\cfoot{\thepage}

% Personalizzazione intestazione
\makeatletter
\renewcommand{\chaptermark}[1]{%
  % Mostra solo "Numero. Titolo Capitolo" (es: 1. Contesto aziendale) ed elimina "Capitolo"
  \markboth{\ifnum\c@secnumdepth>\m@ne \if@mainmatter \thechapter.\ \fi\fi #1}{}%
}
\renewcommand{\sectionmark}[1]{%
  % Mostra solo il nome della sezione, eliminando completamente i numeri (es: Wintech Spa)
  \markright{#1}%
}
\makeatother


% Custom hyphenation rules
\hyphenation {
    e-sem-pio
    ex-am-ple
}

% Images path, not using \graphicspath because it doesn't properly work with
% latexmk custom dependencies
\NewCommandCopy{\latexincludegraphics}{\includegraphics}
\renewcommand{\includegraphics}[2][]{\latexincludegraphics[#1]{../images/#2}}

% Page format settings
\setlength{\parindent}{14pt}    % first row indentation
\setlength{\parskip}{0pt}       % paragraphs spacing


% Load variables
\newcommand{\myName}{Giulia Barzon}
\newcommand{\myID}{2101074}
\newcommand{\myTitle}{Data Governance e Sviluppo Sicuro: Integrazione di OpenMetadata nell'Infrastruttura Aziendale}
\newcommand{\myDegree}{Tesi di Laurea}
\newcommand{\myUni}{Università degli Studi di Padova}
% For BSc level just use "Corso di Laurea" and don't add "Triennale" to it
\newcommand{\myFaculty}{Corso di Laurea in Informatica}
\newcommand{\myDepartment}{Dipartimento di Matematica ``Tullio Levi-Civita''}
\newcommand{\profTitle}{Prof.}
\newcommand{\myProf}{Michele Scquizzato}
\newcommand{\myLocation}{Padova}
\newcommand{\myAA}{2025-2026}
\newcommand{\myTime}{Settembre 2026}

% PDF file metadata fields
% when updating them delete the build directory, otherwise they won't change
\begin{filecontents*}{\jobname.xmpdata}
  \Title{Data Governance e Sviluppo Sicuro: Integrazione di OpenMetadata nell'Infrastruttura Aziendale}
  \Author{Giulia Barzon}
  \Language{it-IT}
  \Subject{Analisi, valutazione e integrazione del Data Catalog OpenMetadata nell'infrastruttura dell'azienda Wintech S.p.A. Il progetto si focalizza sull'automazione della Data Governance, il rispetto dei principi dello Sviluppo Sicuro per la tutela dei dati sensibili e personali e l'integrazione con la Business Intelligence.}
  \Keywords{OpenMetadata\sep Data Governance\sep Sviluppo Sicuro\sep Data Catalog\sep Data Quality\sep Dati Sensibili\sep GDPR\sep Data Observability\sep Business Intelligence}
\end{filecontents*}


% % ACRONIMI SEMPLICI
\newacronym{bi}{BI}{Business Intelligence}
\newacronym{samm}{SAMM}{Software Assurance Maturity Model}


% TERMINI CON ACRONIMO + GLOSSARIO
% GDPR
\newacronym[description={\glslink{gdprg}{General Data Protection Regulation}}]
    {gdpr}{GDPR}{General Data Protection Regulation}

\newglossaryentry{gdprg} {
    name=\glslink{gdpr}{GDPR},
    text=General Data Protection Regulation,
    sort=gdpr,
    description={Regolamento generale sulla protezione dei dati personali emanato dall'Unione Europea (Regolamento UE 2016/679). Stabilisce le norme relative alla raccolta, al trattamento, alla conservazione e alla cancellazione dei dati personali, imponendo alle organizzazioni precisi obblighi in materia di trasparenza, consenso e sicurezza dei dati.}
}

% PII
\newacronym[description={\glslink{piig}{Personally Identifiable Information}}]
    {pii}{PII}{Personally Identifiable Information}

\newglossaryentry{piig} {
    name=\glslink{pii}{PII},
    text=Personally Identifiable Information,
    sort=pii,
    description={Qualsiasi dato che consenta di identificare direttamente o indirettamente una persona fisica (es. nome, cognome, indirizzo e-mail, numero di telefono, codice fiscale). Nel contesto del Data Catalog, i dati PII richiedono una gestione particolare in conformità alle normative vigenti.}
}

% RBAC
\newacronym[description={\glslink{rbacg}{Role-Based Access Control}}]
    {rbac}{RBAC}{Role-Based Access Control}

\newglossaryentry{rbacg} {
    name=\glslink{rbac}{RBAC},
    text=Role-Based Access Control,
    sort=rbac,
    description={Modello di controllo degli accessi che assegna i permessi agli utenti in base al ruolo che ricoprono nell'organizzazione, anziché individualmente. In piattaforme come OpenMetadata, il sistema RBAC determina quali operazioni possono essere eseguite su quali risorse.}
}


% VOCI DI GLOSSARIO
\newglossaryentry{anonimizzazione}{
    name={Anonimizzazione},
    description={Processo irreversibile mediante il quale i dati personali vengono alterati in modo tale da non consentire l'identificazione, diretta o indiretta, della persona fisica a cui si riferiscono. Una volta anonimizzati, i dati non sono più soggetti alle disposizioni del GDPR.}
}

\newglossaryentry{pseudonimizzazione}{
    name={Pseudonimizzazione},
    description={Processo mediante il quale i dati personali vengono trattati in modo tale da non poter essere attribuiti a una persona fisica specifica senza l'utilizzo di informazioni aggiuntive. A differenza dell'anonimizzazione, la pseudonimizzazione è un processo reversibile e i dati rimangono soggetti al GDPR.}
}

\newglossaryentry{dataquality}{
    name={Data Quality},
    description={Misura che valuta quanto un insieme di dati sia "sano" e adatto allo scopo per cui deve essere utilizzato dall'azienda. Applicare la Data Quality significa adottare un sistema che implementa controlli automatici sui dati per individuare anomalie, dati incompleti o informazioni incoerenti.}
}

\newglossaryentry{datagovernance}{
    name={Data Governance},
    description={Insieme di policy, ruoli, regole e processi aziendali che garantiscono che i dati siano sicuri, affidabili, accessibili e conformi alle normative. Rappresenta il "regolamento aziendale" che gestisce il ciclo di vita del dato.}
}

\newglossaryentry{dataingestion}{
    name={Data Ingestion},
    description={Processo di estrazione e importazione automatica dei metadati dalle fonti dati collegate al catalogo. Include la raccolta di informazioni strutturali (tabelle, colonne) ed eventuale profilazione e classificazione dei dati.}
}

\newglossaryentry{datalineage}{
    name={Data Lineage},
    description={Rappresenta la tracciabilità del dato lungo il suo intero ciclo di vita: indica l'origine del dato, le trasformazioni che ha subito e i sistemi attraverso cui è transitato prima di giungere alla destinazione finale, aumentando la fiducia nell'affidabilità delle informazioni.}
}

\newglossaryentry{dataobservability}{
    name={Data Observability},
    description={Capacità di conoscere in tempo reale lo "stato di salute" dei dati. Applica i concetti di monitoraggio del software ai flussi di dati, permettendo di identificare, diagnosticare e risolvere i problemi prima che impattino gli utenti aziendali.}
}

\newglossaryentry{datadiscovery}{
    name={Data Discovery},
    description={Processo di ricerca e centralizzazione volto a scoprire quali informazioni sono presenti all'interno dei vari sistemi dati aziendali. In una piattaforma di Data Catalog, agisce come un motore di ricerca intuitivo per esplorare i dataset disponibili.}
}

\makeglossaries
\defglsentryfmt[main]{\glsgenentryfmt\ifglsused{\glslabel}{}{\glsfirstoccur}}
\loadglsentries{appendix/glossary-entries}
% ESECUZIONE FORZATA NEL CASO NON ANDASSE IL GLOSSARIO: makeglossaries -d build thesis

\bibliography{appendix/bibliography}

\defbibheading{bibliography} {
    \cleardoublepage
    \phantomsection
    \addcontentsline{toc}{chapter}{\bibname}
    \chapter*{\bibname\markboth{\bibname}{\bibname}}
}

% Spacing between entries
\setlength\bibitemsep{1.5\itemsep}

% CATEOGRIE E VOCI DELLA BIBLIOGRAFIA
\DeclareBibliographyCategory{opere}
\DeclareBibliographyCategory{web}
\DeclareBibliographyCategory{report}

\addtocategory{web}{site:wintech}
\addtocategory{web}{site:agile-manifesto}
\addtocategory{report}{gartner:data-catalog-2017}


\defbibheading{opere}{\section*{Riferimenti bibliografici}}
\defbibheading{web}{\section*{Siti Web consultati}}
\defbibheading{report}{\section*{Report e studi}}


\captionsetup{
    tableposition=top,
    figureposition=bottom,
    font=small,
    format=hang,
    labelfont=bf
}

\hypersetup{
    %hyperfootnotes=false,
    %pdfpagelabels,
    colorlinks=true,
    linktocpage=true,
    pdfstartpage=1,
    pdfstartview=,
    breaklinks=true,
    pdfpagemode=UseNone,
    pageanchor=true,
    pdfpagemode=UseOutlines,
    plainpages=false,
    bookmarksnumbered,
    bookmarksopen=true,
    bookmarksopenlevel=1,
    hypertexnames=true,
    pdfhighlight=/O,
    %nesting=true,
    %frenchlinks,
    urlcolor=webbrown,
    linkcolor=RoyalBlue,
    citecolor=webgreen
    %pagecolor=RoyalBlue,
}

% Delete all links and show them in black
\if \isprintable 1
    \hypersetup{draft}
\fi

% Listings setup
\lstset{
    language=[LaTeX]Tex,%C++,
    keywordstyle=\color{RoyalBlue}, %\bfseries,
    basicstyle=\small\ttfamily,
    %identifierstyle=\color{NavyBlue},
    commentstyle=\color{Green}\ttfamily,
    stringstyle=\rmfamily,
    numbers=none, %left,%
    numberstyle=\scriptsize, %\tiny
    stepnumber=5,
    numbersep=8pt,
    showstringspaces=false,
    breaklines=true,
    frameround=ftff,
    frame=single
}

\definecolor{webgreen}{rgb}{0,.5,0}
\definecolor{webbrown}{rgb}{.6,0,0}

\newcommand{\sectionname}{sezione}
\addto\captionsitalian{\renewcommand{\figurename}{Figura}
                       \renewcommand{\tablename}{Tabella}}

\newcommand{\glsfirstoccur}{\ap{{[g]}}}

% ORIGINALE: \newcommand{\intro}[1]{\emph{\textsf{#1}}}
\newcommand{\intro}[1]{\emph{#1}} % SENZA SANS-SERIF

% Risks environment
\newcounter{riskcounter}                % define a counter
\setcounter{riskcounter}{0}             % set the counter to some initial value

%%%% Parameters
% #1: Title
\newenvironment{risk}[1]{
    \refstepcounter{riskcounter}        % increment counter
    \par \noindent                      % start new paragraph
    \textbf{\arabic{riskcounter}. #1}   % display the title before the content of the environment is displayed
}{
    \par\medskip
}

\newcommand{\riskname}{Rischio}

\newcommand{\riskdescription}[1]{\textbf{\\Descrizione:} #1.}

\newcommand{\risksolution}[1]{\textbf{\\Soluzione:} #1.}

% Use case environment
\newcounter{usecasecounter}             % define a counter
\setcounter{usecasecounter}{0}          % set the counter to some initial value

%%%% Parameters
% #1: ID
% #2: Nome
\newenvironment{usecase}[2]{
    \renewcommand{\theusecasecounter}{\usecasename #1}  % this is where the display of
                                                        % the counter is overwritten/modified
    \refstepcounter{usecasecounter}             % increment counter
    \vspace{10pt}
    \par \noindent                              % start new paragraph
    {\large \textbf{\usecasename #1: #2}}       % display the title before the
                                                % content of the environment is displayed
    \medskip
}{
    \medskip
}

\newcommand{\usecasename}{UC-}

\newcommand{\usecaseactors}[1]{\textbf{\\Attore Principale:} #1. \vspace{4pt}}
\newcommand{\usecasedesc}[1]{\textbf{\\Descrizione:} #1. \vspace{4pt}}
\newcommand{\usecasepre}[1]{\textbf{\\Precondizioni:} #1. \vspace{4pt}}
\newcommand{\usecasepost}[1]{\textbf{\\Postcondizioni:} #1. \vspace{4pt}}
\newcommand{\usecasemain}[1]{\textbf{\\Scenario principale:} \vspace{2pt} #1}
\newcommand{\usecasealt}[1]{\textbf{\\Scenario alternativo:} \vspace{2pt} #1}

% Namespace description environment
\newenvironment{namespacedesc}{
    \vspace{10pt}
    \par \noindent  % start new paragraph
    \begin{description}
}{
    \end{description}
    \medskip
}

\newcommand{\classdesc}[2]{\item[\textbf{#1:}] #2}
